The \href{https://www.st.com/resource/en/datasheet/stm32n657a0.pdf}{STM datasheet} and \href{https://www.st.com/resource/en/application_note/an5967-getting-started-with-hardware-development-for-stm32n6-mcus-stmicroelectronics.pdf}{getting started handbook} were used to identify pins and their purpose. Most circuits and component values were obtained from the guide or datasheet.
For example, pg. 7 of the getting started guide details the decoupling power figures required on most of the supply pins. For pins that were not explicitly described in the guide, I used the datasheet to find the pins. The decoupling inductor on VDDAON (near pin E1 on the schematic) had no specified value, so, I copied the inductor value used on the VDDCORE pull-down set.
The schematic symbol was split in two to make it easier to connect the GPIO pins (opposed to the BGA order it was previously in: A1, A2, ..., B1, B2, ..., which made it difficult to cleanly connect components).
\nl
Additionally, debug LEDs were added to the schematic to make it easier to know which lines are active and which aren't. If I test one of the power lines once and it provides us with the correct voltage, then whenever the corresponding LED is active, I know it has the correct voltage.
Similarly, an LED was added for reading from the sensor and writing to the uSD card. Resistor arrays were used rather than individual resistors to save on the number of components on the PCBs, but MOSFET arrays weren't used because their biasing voltage is typically higher than the regular MOSFET bias voltage (~0.7V), so lines like 0.8V would not bias the transistor and activate the LED.
Each LED will use the HSMH-C190 (green). According to the datasheet, the maximum current rating is 25mA, so with a supply voltage of 3V3, I require a minimum of 68$\Omega$ with the MOSFET configuration. I selected values of 470$\Omega$ since it's a very common resistor value and it satisfies the requirements without being too bright (too bright and they'd begin to blind us after a while).
Although the MOSFET set will be slightly dimmer than the other, their brightness difference is not a concern.
A high value pull-down resistor was included on the base of the MOSFETs to ensure that the MOSFET is not active when the driving signal is not actively pulling the gate to a defined voltage (also used a resistor array). Since the MOSFETs are used on non-time-critical areas (like debug LEDs and switches), the time difference between $10\texttt{k}\Omega$ and lower resistances was not a concern. I picked $10\texttt{k}\Omega$ since it is a common resistor value.
\nl
The 3V3, 1V8 and VDDAON enable MOSFETs use a different package since their drain current exceeds the maximum of the DMG1012T. Instead, I have selected the \href{https://au.mouser.com/ProductDetail/Infineon-Technologies/BSS214NWH6327XTSA1?qs=Dj8dEyEphkMe%252Bfnz0ZE2og%3D%3D}{BSS214NWH6327XTSA1} for the power domains (continuous drain current of up to 1.5A)